\documentclass[11pt]{article}
\usepackage{amsmath,amsfonts}
\usepackage{graphicx}
%\usepackage{xcolor}
%\definecolor{gray}{rgb}{.75,.75,.75}
%\usepackage[landscape]{geometry}
%\usepackage{hyperref}
%\usepackage[bookmarks=true, pdfborder={0 0 .5}, urlbordercolor=gray]{hyperref}

\renewcommand{\thefootnote}{\fnsymbol{footnote}}

\addtolength{\topmargin}{-.5in}
\addtolength{\textheight}{.9in}
\addtolength{\oddsidemargin}{-.3in}
\addtolength{\textwidth}{.6in}
\pagestyle{empty}

\newcommand{\ds}{\displaystyle}
\newcommand{\ts}{\textstyle}

\newcommand{\Z}{\mathbb{Z}}
\newcommand{\R}{\mathbb{R}}
\newcommand{\C}{\mathbb{C}}
\newcommand{\EE}{\mathcal{E}}
\newcommand{\tZ}{\tilde{\Z}}
\newcommand{\ep}{\epsilon}

\newcommand{\sech}{\mathrm{sech}}
\newcommand{\csch}{\mathrm{csch}}


\newcommand{\Sym}{\mathrm{Sym}}

\renewcommand{\circle}[1]{{\bigcirc\hspace{-.75em}#1}}

\begin{document}

\footnote[0]{Notes by Peter Magyar \ {\tt magyar@math.msu.edu}}
\vspace{-1em}

\centerline{\bf Math 133\hspace{1.5em} \hfill{\Large\bf
More Power Series}\hfill Stewart \S11.9}
\vspace{1em}

\noindent 
{\bf Calculus on power series.}
In this section, we will finally be able to give series for some interesting numbers and functions.  The cornerstone we build on is the geometric series $\frac1{1-r}=\sum_{n=0}^\infty r^n$, which we manipulate into much more interesting series formulas. Now, series of {\it numbers} can only be manipulated by algebra; 
but we have introduced series of {\it functions} (power series), 
where we can apply the calculus operations of differentiation and integration.

\begin{quote}
{\it Theorem:} Given a power series convergent for 
$|x{-}a|<R$, for some $R>0$:
$$
f(x) 
\ =\
\sum_{n=0}^\infty c_n (x{-}a)^n
\ =\
c_0 + c_1(x{-}a) + c_2(x{-}a)^2+ c_3(x{-}a)^3 +\cdots .
$$
Then we have:
\begin{itemize}
\item For $|x{-}a|<R$, the derivative is:
$$
f'(x) 
\ =\
\sum_{n=1}^\infty n\,c_n (x{-}a)^{n-1}
\ =\
c_1 + 2c_2(x{-}a)+ 3c_3(x{-}a)^2 +\cdots .
$$

\item For $|x{-}a|<R$, the antiderivative is:
$$
\int\!\! f(x) \, dx 
\ =\
C \ +\ 
\sum_{n=0}^\infty c_n \frac{(x{-}a)^{n+1}}{n{+}1\ \ }
%$$
%$$
\ = \ C\ +\  c_0 (x{-}a) \ +\  c_1\frac{(x{-}a)^2}{2 \ }\ +\  c_2\frac{(x{-}a)^3}{3 \ } \ +\ \cdots .
$$
\end{itemize}

\end{quote}

\noindent
That is, we can differentiate and integrate a power series term by term.
\\[1em]
{\bf Derivatives of geometric series.}
We think of the geometric series as a function:
$$
f(x) \ =\ \frac{1}{1{-}x} \ =\ \sum_{n=0}^\infty x^n \ = \ 1 + x + x^2 + x^3 +x^4\cdots,
$$
Its first two derivatives are, for $|x|<1$:
$$
f'(x) \ =\ \frac{1}{(1{-}x)^2} \ =\ \sum_{n=0}^\infty nx^{n-1} \ = \ 1 + 2x + 3x^2 + 4x^3 +\cdots,
$$
$$
f''(x) \ =\ \frac{2}{(1{-}x)^3} \ =\ \sum_{n=0}^\infty n(n{-}1)x^{n-2} \ = \ 2 + 3\cdot2x + 4\cdot3x^2 +\cdots.
$$
This lets us find the sums of many series similar to geometric series.
\\[.5em]
{\sc example:} Find the sum $\ds\sum_{n=1}^\infty \frac{n^2}{3^n} = \frac{1}{3} + \frac{4}{9} + \frac{9}{27} + \frac{16}{81}+\cdots$. The first 10 terms sum to about $1.499$, so we can guess the answer is $\frac32$, but how to be sure?
Consider the series as one output of a function;
it is $g(\frac13)$ for:
$$
g(x) \ = \ \sum_{n=0}^\infty n^2x^n  \ =\ x+ 4x^2 + 9 x^3+ 16x^4+\cdots.
$$
We can write this in terms of our known series because $n^2 = n(n{-}1) + n$, so that:
$$\begin{array}{rcl}
g(x) \ = \ \sum_{n=0}^\infty n^2x^n &=&
 \sum_{n=0}^\infty n(n{-}1)x^n \ +\ \sum_{n=0}^\infty nx^n
\\[.5em]
 &=&
 x^2\,\sum_{n=2}^\infty n(n{-}1)x^{n-2} \ +\ x\,\sum_{n=1}^\infty nx^{n-1}
\\[.5em]
 &=&
 x^2\, \frac{2}{(1{-}x)^3} + x\, \frac{1}{(1{-}x)^2} 
\ =\ 
\frac{x(x{+}1)}{\ (1{-}x)^3}\,.
\end{array}$$
Hence our series sums to $g(\frac13) = \frac{\frac13(\frac13{+}1)}{\ (1{-}\frac13)^3}=\frac32$.  Fun!
\\[1em]
{\bf Integrals of geometric series.} Using the second part of the Theorem, we integrate
the geometric series formula $\frac 1{1-x}=\sum_{n=0}^\infty x^n$:
$$
\int\! \frac1{1-x}\,dx \ =\ -\log(1{-}x) \ =\ C+ \sum_{n=0}^\infty \frac1{n{+}1} x^{n+1} \ = \ C + x + \frac{x^2}2 + \frac{x^3}3 +\frac{x^4}4\cdots.
$$
To determine the correct constant $C$, we use the initial value at $x=0$:
the left side becomes $-\log(1-0)=0$, the right side $C + 0 +\frac{\,0^2}{2}+\cdots=C$, so $C=0$. We can manipulate this into an expression for $\ln(x)$ itself:
$$\begin{array}{rcl}
\ln(1{+}x) \ =\ 
-\left(-\ln(1{-}(-x))\right)
&=& 
-\left((-x) +\frac{(-x)^2}{2} +\frac{(-x)^3}{3} +\frac{(-x)^4}{4}+\cdots \right)
\\[.5em]
&=& 
x -\frac12 x^2 +\frac13 x^3 -\frac14 x^4 +\cdots 
\\[.5em]
\ln(x) \ =\ 
\ln(1{+}(x{-}1)) 
&=&
(x{-}1) -\frac12(x{-}1)^2 +\frac13 (x{-}1)^3 -\frac14 (x{-}1)^4 +\cdots
\end{array}$$
The first series converges for $|x|<1$, 
so plugging $x{-}1$ in place of $x$,
we find that the last series converges for $|x{-}1|<1$.
We conclude:
$$
\ln(x) \ =\ \sum_{n=1}^\infty \frac{(-1)^n}{n}(x{-}1)^n \quad\text{for}\quad |x{-}1|<1.
$$
For example, taking $x=\frac12$, we get $x{-}1=-\frac12$, and:
$$
\ln(\tfrac12)\ =\  (-\tfrac12) - \tfrac12(-\tfrac12)^2 + \tfrac13(-\tfrac12)^3 -
\tfrac14(-\tfrac12)^4+\cdots.
$$
Since $\ln(\frac12)=-\ln(2)$, we conclude:
$$
\ln(2) \ =\ 
\sum_{n=1}^\infty \frac1{n\,2^n}
\ =\
\frac12 +\frac1{2\cdot2^2}+\frac1{3\cdot2^3}+\frac1{4\cdot2^4}+\cdots.
$$
Since the denominator grows quickly, the series converges rapidly and gives a good approximation to $\ln(2)=0.69314\cdots$ after only a few terms.
For example, the first 10 terms give:
$\sum_{n=1}^{10} \frac1{n\,2^n}=0.69306\cdots$, accurate to 3 or 4 decimal places.
When a calculator or Wolfram Alpha computes logarithms, it is using some method similar to this, taking enough terms to obtain the desired number of decimal places.
\newpage

\noindent
{\bf Series for $\boldsymbol{\pi}$.}  We obtained the series for logarithm because it is an inverse function whose derivative is a rational function (see the Inverse Derivative Theorem in \S6.1). We can do the same for the arctangent:
$$
\frac1{1{+}x^2} \ =\ \frac{1}{1{-}(-x)^2}
\ =\ \sum_{n=0}^\infty (-x^2)^n\ =\ \sum_{n=0}^\infty (-1)^nx^{2n}
\ =\ 1 -x^2 + x^4 -x^6 + \cdots
$$
$$
\tan^{-1}(x) \ =\ \int\!\frac1{1{+}x^2}\ dx
\ =\ \sum_{n=0}^\infty(-1)^n\frac{x^{2n+1}}{2n{+}1}
\ =\ x -\frac{x^3}{3} +\frac{x^5}5-\frac{x^7}7+\cdots.
$$
We know there is no constant shift because both sides are 0 at $x=0$. 
The series converges for $|x|<1$.
Since $\tan(\frac\pi4) = 1$, we have:
$$
\frac{\pi}4 \ = \ \tan^{-1}(1) 
\ =\ 1 -\frac{1}{3} +\frac{1}5-\frac{1}7+\cdots.
$$
Here we must be careful, since $x=1$ is all the way at the edge 
of the interval of convergence: the series is not absolutely convergent,
but we can show conditional convergence by the Alternating Series Test (\S11.6/II).
(Also, Abel's Theorem\footnote{
{\it Theorem:} Let $f(x)=\sum_{n=0}^\infty a_n x^n$ converge for $|x|<R$,
and suppose $f(R)=\sum_{n=0}^\infty a_nR^n$ exists. Then $\lim_{x\to R^-} f(x) = f(R)$,
so the power series defines a left-continuous function at $x=R$.
\\
{\it Proof:} Replacing $f(x)$ with $f(Rx)$, we can assume $R=1$. Replacing $a_0$ with
$a_0-\sum_{n=1}^\infty a_n$, we can assume $f(1)=\sum_{n=0}^{\infty} a_n = 0$. 
We must show $\lim_{x\to 1^-} f(x) \stackrel{\text{def}}=
\lim_{x\to 1^-} \sum_{n=0}^\infty a_nx^n \stackrel?= 0$.

Abel's summation-by-parts formula: for sequences $\{a_n\}$, 
$\{b_n\}$, with $s_n=\sum_{k=0}^n a_k$, we have
 $\sum_{n=0}^N a_nb_n 
= s_N b_{N+1}  -  \sum_{n=0}^{N} s_n(b_{n+1}{-}b_n)$. 
For $b_n=x^n$ with $b_{n+1}-b_n = (x{-}1)x^n$, we have
$$\ts
\sum_{n=0}^N a_nx^n  
\ = \ s_Nx^{N+1}  +(1{-}x)\sum_{n=1}^{N} s_{n}x^n
.$$
For fixed $0<x<1$,  $N\to\infty$,
$f(x) = (1{-}x)\sum_{n=0}^{\infty} s_{n}x^n$.
Given any $\ep>0$, we must choose $x$ close enough to 1 to force 
$|f(x)| = (1{-}x) \left|\sum_{n=0}^{\infty} s_{n}x^n\right| < \ep$.
First choose $N$ large enough that 
$|s_n|=\left|\sum_{k=0}^n a_n\right|<\frac12\ep$ for $n\geq N$.
Then choose $(1{-}x)$ small enough that 
\smash{$(1{-}x)\left|\sum_{n=0}^{N-1} s_{n}x^n\right|<\frac12\ep$:}
$$\begin{array}{rcl}
|f(x)|&\leq& (1{-}x)\,\left|\smash{\sum_{n=0}^{N-1}} s_{n}x^n\right| +
(1{-}x)\, |\sum_{n=N}^{\infty} s_{n}x^n|
\\[.5em]
&\leq& \frac12\ep + (1{-}x)\,\sum_{n=N}^\infty\frac12\ep\, x^n
\ =\ \frac12\ep + \frac12\ep(1{-}x) \frac{x^N}{1-x} \ <\ \ep. \qquad\square
\end{array}$$
Now, for $f(x) = x-\frac{x^3}{3}+\frac{x^5}{5}-\cdots$, 
we know that $f(x)=\tan^{-1}(x)$ for $|x|<1$, and Abel's Theorem
gives $f(1) = \lim_{x\to1^-} f(x) = \lim_{x\to1^-} \tan^{-1}(x) = \tan^{-1}(1) =\frac\pi 4$.
}
shows that the convergent series gives the expected value $\tan^{-1}(1)$.)

This is known as the Leibnitz formula. It is astonishing because the series is so simple and 
seemingly has no relation to circles or angles.
It can be used to compute $\pi$ (multiplying both sides by 4),
but it is inefficient because of the slow convergence. It takes about 200 terms to get 2 decimal places of accuracy, $\pi\approx 3.14$. 

Machin's trick gives an efficient series.
Angle addition $\tan(\alpha+\beta) = \frac{\tan\alpha \, +\, \tan\beta}{1-\tan\alpha\tan\beta}$
for $x=\tan\alpha$, $y=\tan\beta$ gives: 
$\ts \arctan x + \arctan y = \arctan(\frac{x+y}{1-xy}).$
We want small $x,y$ with $\frac{x+y}{1-xy} = 1$, to get $\arctan x + \arctan y = \arctan 1 = \frac\pi4$.
This is the hyperbola $(x{+}1)(y{+}1) = 2$, with a rational point $(x,y) = (\frac12,\frac13)$. Thus
$$\ts
\pi = 4\left(\arctan(\frac12) + \arctan(\frac13)\right) = \sum\limits_{n=0}^\infty (-1)^n \frac{4}{2n+1}(\frac1{2^{2n+1}} + \frac1{3^{2n+1}}).
$$
Up to $n=3$, this gives $\pi$ to 2 decimal places; $n=15$ gives 10 places:
3.1415926535\ldots

\end{document}

%%%%%%%%%%%  Picture  %%%%%%%%%%
\\[0em]
\centerline{
%\includegraphics[width=2in]
{1-8-Discont-iv.png}
\qquad \qquad 
%\includegraphics[width=2in]
{1-8-Discont-v.png} 
}
\vspace{0em}

\noindent
%%%%%%%%%%%%%%%%%%%%%%%%%

%%%%%%%%%%%  Table  %%%%%%%%%%
In fact, we get the following derivatives:
$$\begin{array}{|c||c|c|c|c|c|c|}
\hline &&&&&& \\[-.9em] 
f(x) & \sin(x) & \cos(x) & \tan(x)  & \sec(x) & \csc(x) & \cot(x)
\\[.2em] \hline &&&&&& \\[-.9em]
f'(x) &\cos(x) & -\sin(x) & \sec^2(x) & \tan(x)\sec(x) & -\cot(x)\csc(x) & -\csc^2(x)
\\[.1em] \hline
\end{array}$$
\\[1em]
%%%%%%%%%%%%%%%%%%%%%%%%%%



















